\documentclass{article}

\usepackage[T1]{fontenc}
\usepackage{xcolor}
\usepackage{microtype}
\usepackage[colorlinks=true,linkcolor=blue,citecolor=teal]{hyperref}

\title{A Small Study of Reproducible Typesetting}
\author{Ada Example}
\date{June 2026}

\begin{document}

\maketitle

\begin{abstract}
This note demonstrates a compact LaTeX workflow for comparing manuscript
versions.  The old version argues that reproducible typesetting is mainly a
matter of keeping source files in one directory and compiling them in a stable
order.
\end{abstract}

\section{Motivation}

Scientific documents often combine prose, figures, and bibliographic records.
When these ingredients are revised independently, reviewers need a clear view
of what changed.  Tools such as \texttt{latexdiff} are useful because they mark
insertions and deletions directly in the compiled document.

The baseline workflow in this example is intentionally simple.  Authors edit a
single \texttt{.tex} file, keep a generated \texttt{.bbl} file next to it, and
compile the document with \texttt{latexmk}.  Classic references on literate
typesetting and LaTeX practice motivate this approach \cite{knuth1984,lamport1994}.

\section{Method}

The document is built from deterministic inputs.  The bibliography is stored as
a checked-in \texttt{.bbl} file so that the example does not require running
BibTeX or Biber.  This makes the comparison portable across small test
environments.

\section{Conclusion}

The old workflow is adequate for short examples, but it leaves room for a more
explicit build rule that treats the bibliography as a first-class dependency.

\input{\jobname.bbl}

\end{document}
